\section{Outlook}

\green{Hi Caroline, in case you're already reading this, I finally got the CNN to train on the 6-param inference problem as well! I can now pretrain with a smaller bottlenck too and I have my first successful training runs, I'm still implementing that in the thesis and I'm excited to see the results tonight (Sat.) $\smiley$} 

See~\ref{fig:joint_loss_example} for an example of joint training loss after pretraining both models on a bottleneck dimension of 5.

\begin{figure}
    \centering
    \includegraphics[width=\linewidth]{graphics/cnn_5/first_run.png}
    \caption{Joint loss for bottleneck dim 5}
    \label{fig:joint_loss_example}
\end{figure}


How one could incorporate information geometry for similar problems

Any idea how to numerically resolve $F_{ij}$ for this flow?

Alternative choices than to model instead of $p(x|y)$ and $p(x)$?

The compressing network usually has a much heavier architecture. Once it has been trained, one can learn different lighter networks to map its summary to different target distributions. \red{This is only really true in the self-supervised case, see SKATR. So Maybe this should go in outlook}